\subsection{An Euler precheck using only quadrilateral rays}
\label{sec:euler-corners}

Complete standard-ray enumeration is more work than a query asking only
whether a non-link ray has positive Euler characteristic. The following
elementary consequence of canonical extension separates these two tasks.
Its ingredients are classical: canonical surfaces have no vertex-link
components, and canonical quadrilateral vertex surfaces are standard vertex
surfaces \cite[Section 3 and Lemma 4.5]{burtonconversion}. The explicit
convexity formulation below is used as a proposed optimization; no claim
of literature priority is made.

\begin{theorem}[Objectives nonnegative on links attain a maximum at Q-corners]
\label{thm:euler-corners}
Let $\lambda$ be a linear functional on standard matching vectors such
that $\lambda(L_v)\ge0$ at every global vertex. On the supplied sector
cone, $F_\lambda(q)=\lambda(\lift(q))$ is positively homogeneous and
convex. If the normalized section $P$ is nonempty, then
\[
 \max_{q\in P}F_\lambda(q)
   =\max_{a\in\operatorname{Vert}(P)}F_\lambda(a).
\]
Every corner lift $\lift(a)$ spans a non-link standard extreme ray.
For a compact finite three-manifold triangulation, this applies to Euler
characteristic, because each vertex link is a disc or sphere. Consequently
there is a non-link standard ray of positive Euler characteristic if and
only if some canonical lift of a quadrilateral-cone extreme ray has
positive Euler characteristic. The convexity and equivalence hold in
all matching dimensions.
\end{theorem}

\begin{proof}
Let $U(q)$ be the signed linear matching vector with triangle entries
$\ell_c(q)$ and quadrilateral entries $q$, with zero offsets in every
vertex group. Include the singleton groups with potential zero. Then
\[
 \lift(q)=U(q)-\sum_v m_v(q)L_v,
 \qquad
 F_\lambda(q)=\lambda(U(q))-
       \sum_v\lambda(L_v)\min_{c\in G_v}\ell_c(q).
\]
Each minimum of linear functions is concave and positively homogeneous.
The coefficients $\lambda(L_v)$ are nonnegative, so the displayed
expression is convex and positively homogeneous. If
$q=\sum_i\alpha_i a_i$ is a convex combination of corners of $P$, then
\[
 F_\lambda(q)\le\sum_i\alpha_i F_\lambda(a_i)
              \le\max_i F_\lambda(a_i).
\]
The reverse maximum inequality follows since every corner lies in $P$.

A corner of $P$ spans an extreme ray of $\cone$. Choose any anchor cell
containing it. That direction remains extreme in the anchor cell, so
Theorem~\ref{thm:face} makes its lift standard-extreme. Conversely, every
non-link standard extreme ray is canonical. Euler characteristic is a
homogeneous linear functional in standard normal coordinates
\cite[Section 2 and Lemma 10]{burtonozlen}, and its values on the links
are one or two. Applying the maximum equality and positive rescaling
proves the claimed positive-Euler equivalence.
\end{proof}

The comparison concerns the normalized quantity
$\chi(\lift(q))/(\mathbf1^Tq)$ when primitive integer directions are used.
Primitive normalization can change the value of Euler characteristic by
a positive scale; it does not change its sign. Unnormalized numerical
maxima of primitive representatives must not be substituted for the
maximum in the theorem.

\begin{corollary}[A negative Euler answer excludes essential normal discs]
\label{cor:no-disc-corners}
If all canonical corner lifts have nonpositive Euler characteristic,
there is no essential normal disc anywhere in the supplied sector.
\end{corollary}

\begin{proof}
An essential normal disc is connected and is not a vertex-link disc:
the latter has boundary bounding the small disc about a boundary vertex.
Hence it has no vertex-link component and is canonical
\cite[Definitions 3.1 and 3.9, Lemma 3.5]{burtonconversion}. Its
nonzero quadrilateral vector would give $F_\chi>0$, contradicting
Theorem~\ref{thm:euler-corners}. The empty-section case has only
vertex-link surfaces and gives the same conclusion.
\end{proof}

A positive answer has a different meaning. The returned standard ray
may be a sphere, a projective plane, or an inessential disc. Its essentiality
still requires the compressed topological observer. The theorem does
not assert that an essential disc must itself occur among the canonical
quadrilateral corner lifts. It also does not claim that all standard
rays occur there; the additional overlay vertices remain necessary for
complete enumeration.

\begin{theorem}[Quadratic post-kernel Euler precheck]
\label{thm:euler-complexity}
For matching nullity at most three, the positive-Euler non-link ray query
can be decided after kernel construction with
\[
 O(1+t+k^2)
\]
rational arithmetic operations using deterministic array bookkeeping.
A positive answer can include one primitive full-coordinate standard ray
within the same bound. This is an arithmetic bound with polynomial bit
complexity; compressed component testing and independent dense replay
are separate costs.
\end{theorem}

\begin{proof}
First build a linear Euler functional $\chi(z)=e^Tz$ in $O(t)$ operations.
Count one face for each elementary normal disc, count normal arcs on one
representative of each physical triangulation face, and count intersection
points on one representative local edge for each global triangulation
edge. Matching makes the representative choices consistent. The formula
$F-E+V$ yields integer coefficients between $-3$ and $5$: a normal disc
meets at most four faces and four edges. The face and edge classes are
obtained by traversing linear-size incidence graphs. This is the usual
bounded-coefficient Euler construction \cite[Lemma 10]{burtonozlen}.

Construct the normalized section by the bounded-chart clipping procedure,
without constructing any minimum diagram. This costs $O(k^2)$ and gives
at most $k$ corners; the empty, point, and segment cases satisfy the same
bound on their number of corners. Preproject the $p=O(k)$ retained
potentials to the at-most-two-dimensional chart in $O(k^2)$ operations.
One pass through the original triangle corners and allowed quadrilateral
entries then aggregates the affine coefficients of $\chi(U(q))$ and the
link coefficients $\chi(L_v)$ in $O(t+k)$ operations. These are constant-size
chart records indexed by the already-known class identifiers.

At a section corner, evaluate the affine term and each group's minimum
using $O(p+k)=O(k)$ operations. Evaluating all corners therefore costs
$O(k^2)$. Theorem~\ref{thm:euler-corners} proves that their maximum decides
the query. If a positive corner is found, its primitive canonical full
lift costs an additional $O(t+k)$ operations. Exact rational chart and
output sizes have the polynomial bit bounds established earlier.
\end{proof}

This Euler precheck is proved and examined by the separate exact-data
audit, but is \emph{not installed in the delivered runtime dispatcher or
certificate APIs}. A production integration should give its negative
answer a certificate checking the complete quadrilateral-corner set and
all Euler values, without pretending that this is a complete standard-ray
list. The measured producer and benchmark source remain unchanged.

